% =============================================================================
% test-abnt.tex
% Documento de teste para abnt.cls (classe base)
%
% Demonstra os comandos públicos da classe base conforme NBR 14724:2024
% e NBR 15287:2025.
%
% Compilação:
%   cd tests && TEXINPUTS=../src: latexmk -pdf -shell-escape -interaction=nonstopmode test-abnt.tex
%
% Checklist:
%   1  — \titulo e \subtitulo preenchem metadados
%   2  — \autor preenche nome do autor
%   3  — \orientador e \coorientador preenchem orientação
%   4  — \instituicao, \local, \ano preenchem dados institucionais
%   5  — \natureza, \programa, \area preenchem natureza do trabalho
%   6  — \imprimircapa gera capa com título em negrito maiúsculas
%   7  — Subtítulo na capa sem negrito, precedido de dois pontos
%   8  — \imprimirfolhaderosto gera folha de rosto com natureza recuada
%   9  — Coorientador aparece na folha de rosto quando definido
%   10 — \pretextual suprime numeração de página
%   11 — \textual restaura numeração de página
%   12 — Cenário sem subtítulo e sem coorientador (segundo par capa+folha)
%   13 — Capa sem subtítulo não exibe dois pontos
%   14 — Folha de rosto sem coorientador não exibe linha de coorientador
% =============================================================================

\documentclass{abnt}

% Metadados — cenário completo (checklist 1–5, 7, 9)
\titulo{Análise de algoritmos de ordenação em grandes volumes de dados}
\subtitulo{um estudo comparativo}
\autor{Maria Clara Santos}
\orientador{Prof. Dr. João Pereira}
\coorientador{Profa. Dra. Ana Beatriz Lima}
\instituicao{Universidade Federal de Minas Gerais}
\local{Belo Horizonte}
\ano{2025}
\natureza{Dissertação apresentada ao Programa de Pós-Graduação em
Ciência da Computação da Universidade Federal de Minas Gerais,
como requisito parcial para obtenção do título de Mestre em
Ciência da Computação.}
\programa{Programa de Pós-Graduação em Ciência da Computação}
\area{Algoritmos e Estruturas de Dados}

\begin{document}

% =====================================================================
% PRÉ-TEXTUAIS (checklist 10)
% =====================================================================
\pretextual

% Capa com subtítulo e coorientador (checklist 6, 7)
\imprimircapa

% Folha de rosto com subtítulo e coorientador (checklist 8, 9)
\imprimirfolhaderosto

% =====================================================================
% CENÁRIO SEM SUBTÍTULO E SEM COORIENTADOR (checklist 12, 13, 14)
% =====================================================================
% Redefinir metadados para cenário mínimo
\titulo{Estudo sobre estruturas de dados probabilísticas}
\subtitulo{}
\coorientador{}

% Segunda capa — sem subtítulo (checklist 13)
\imprimircapa

% Segunda folha de rosto — sem coorientador (checklist 14)
\imprimirfolhaderosto

% =====================================================================
% ELEMENTOS TEXTUAIS (checklist 11)
% =====================================================================
\textual

\chapter{Introdução}

Este trabalho apresenta uma análise comparativa de algoritmos de ordenação
aplicados a grandes volumes de dados. A motivação principal reside na
necessidade de avaliar o desempenho prático em cenários reais.

\chapter{Revisão bibliográfica}

A literatura sobre algoritmos de ordenação é extensa e abrange desde os
métodos clássicos até abordagens mais recentes baseadas em paralelismo.

\chapter{Considerações finais}

Os resultados obtidos confirmam a eficiência assintótica dos algoritmos
analisados e sugerem direções para pesquisas futuras.

\end{document}
